\section{Institutional setting}\label{sec:background}

\subsection{The settlement price}

Equity derivatives on the National Stock Exchange of India expired, during the sample year,
on the last Thursday of each month, or on the preceding trading day when that Thursday was a
holiday. The final settlement price of a stock future and the exercise settlement price of a
stock option are both the closing price of the underlying security in the cash market on the
expiry day. The closing price is not the last trade: it is the volume-weighted average price
of all trades between \DesignWindowStart{} and \DesignWindowEnd{}, the final
\DesignWindowMinutes{} minutes of continuous trading.

Stock derivatives have been settled by delivery since the exchange completed the transition
from cash settlement in 2019. A long futures position held to expiry becomes an obligation to
take delivery of the shares, and an in-the-money stock option becomes a delivery at the
strike. The settlement price therefore determines the final variation margin on every
expiring futures position, and it determines which option series are exercised. For an option
writer whose position sits near a strike, a small movement of the settlement price decides
whether shares must be delivered.

Index derivatives are settled in cash against the closing value of the index, and the
closing value is computed from the closing prices of its constituents. The settlement window
therefore also fixes the payoff of the index options that expire on the same day. In 2022 the
two main index option contracts also expired weekly, on every Thursday. A Thursday that is
not the monthly expiry thus carries the index settlement but not the stock derivatives
settlement. This distinction is used in Section~\ref{sec:design} to separate the two.

\subsection{Why the window may be distorted}

Three features of the arrangement are relevant. The window is known in advance, so a
participant who wishes to trade in it can plan rather than react. It is short, so the volume
of a single participant can be a noticeable share of the total, particularly in a security
that trades little. And because the settlement price is an average rather than a single
print, moving it requires trading throughout the window rather than a single order at the
close. A participant who moves the average does so at a cost, since the price impact of the
trades is incurred on the full quantity while the benefit accrues on the expiring position.

Two responses to the settlement date are not distortions at all. Holders of expiring
positions who do not wish to take or make delivery must close or roll them, and index funds
and arbitrageurs who hold cash positions against futures unwind them at the settlement price
by trading in the window, because that is the price at which their hedge is closed. Both
raise volume in the window, and both are expected to be concentrated in securities with
expiring contracts. The analysis therefore distinguishes between the window being more
active, which these activities predict, and the settlement price being displaced from the
price that prevails before and after, which they do not.

\subsection{Order-level features used}

The cash market is a continuous limit order book with price-time priority. A disclosed
quantity instruction allows a participant to display only part of an order; as the displayed
part is executed it is replenished from the undisclosed remainder, and each replenishment
loses time priority. An immediate-or-cancel instruction executes against the available
depth on arrival and cancels the remainder. Both are recorded in the order file, together
with the category of the submitting participant: custodian, which covers institutional
investors whose trades are cleared through a custodian, proprietary, and other clients
(NCNP).
